\chapter{Medida y doble registro}

\section{Conteo y medida}
\begin{definition}
Si una diferencia pertenece al tránsito, se cuenta; si pertenece a la extensión, se mide. Si mezcla ambas capas, el sistema debe declarar primero el desdoblamiento de lectura.
\end{definition}

\begin{observation}
Esta cláusula separa dos funciones que no deben confundirse: registrar paso y registrar magnitud. La mezcla no queda prohibida, pero sí sometida a declaración explícita.
\end{observation}

\section{Doble registro}
\begin{definition}
La medida estructural \Mmeas\ se expresa mediante un doble registro
\[
\Mmeas := (n,\delta),
\]
donde \(n\) representa el conteo externo y \(\delta\) la subfase interna del proceso de lectura.
\end{definition}

\begin{definition}
Si se requiere una magnitud inducida, se admite la expresión
\[
|\Mmeas| := n + \delta \cdot \eps.
\]
\end{definition}

\begin{observation}
La magnitud inducida no reemplaza al doble registro. Solo proporciona una lectura derivada cuando la comparación de estados lo exige.
\end{observation}

\section{Regla mínima de actualización}
\begin{definition}
La subfase interna \(\delta\) evoluciona dentro del rango discreto declarado por el sistema. Cuando la subfase reinicia su ciclo, el conteo externo \(n\) se incrementa en una unidad.
\end{definition}

\begin{observation}
Esta regla fija una medida con espesor interno. No equivale todavía a una dinámica completa, pero sí establece una forma mínima de acumulación legible.
\end{observation}

\section{Cláusula de alcance}
\begin{observation}
La introducción del doble registro fija una gramática mínima para contar y medir sin colapsar ambas operaciones. No valida por sí sola todos los regímenes posteriores que dependan de medición estructural.
\end{observation}

